\section{Threats to validity}
\label{sec:threats}

\paragraph{Construct validity.}
The central risk is the one we have been explicit about: our labels encode paraphrase and
question-duplication, and we read them as cache-equivalence. The two come apart at the
edges, and the gap runs both ways: a labeled non-paraphrase may occasionally admit the same
answer, which inflates the measured false-hit rate, while a labeled paraphrase may need a
different answer in context, which deflates it. We accept this gap deliberately, because
the alternative, bespoke answer-equivalence annotation at scale, would trade a transparent,
reproducible label for a private one. We therefore do not claim the residual bias favors our
thesis. The claim is narrower: the labels track a public, reproducible notion of
sameness, and the adversarial domain is built to isolate exactly the case a similarity cache
cannot resolve, near-identical wording with divergent meaning. The numbers should be read as
the reliability of the similarity proxy against that human notion of sameness, which is the
proxy a deployed cache actually relies on.

\paragraph{Internal validity.}
Cosine scores are not comparable across encoders, so every cross-encoder statement uses
threshold-free metrics (\prauc{}, \rocauc{}), and every threshold is calibrated per encoder
(\cref{sec:calibration}). Within an encoder, the false-hit rate and coverage are exact
functions of the score distributions, so there is no estimator variance beyond sampling,
which the bootstrap intervals capture.

\paragraph{External validity.}
We measure three domains and six encoders, and the specific numbers will not transfer
verbatim to a new corpus or a next-generation embedding. The contribution that does transfer
is the method: define cache-equivalence, collect a labeled pair set for the target domain,
read the false-hit/coverage frontier, and calibrate the threshold to the measured cost. We
designed the study so that re-running it on a new encoder is a single command. The
adversarial domain is intentionally a worst case rather than an average one; real traffic
will sit between the news and adversarial frontiers, and where it sits is itself something a
deployment can measure with this protocol. Coverage is reported over balanced labeled pairs, so
the absolute hit rate a live system sees will depend on how much of its traffic is genuinely
reusable; the transferable object is the shape of the false-hit/coverage trade-off at a given
ceiling, not the single coverage number.

\paragraph{Conclusion validity.}
All intervals are $95\%$ bias-corrected and accelerated bootstrap intervals over pairs.
Cross-encoder comparisons use a paired bootstrap over the shared pairs and are corrected for
the false discovery rate across the comparison family with the Benjamini--Hochberg procedure;
the lexical-versus-neural gaps on the adversarial domain, for instance, hold at a corrected
$p \le 0.003$. We report the size of each difference next to its test, so that a reliable
difference is not mistaken for a large one. The analysis plan, the primary metric, and the
sample sizes were fixed before the main run and are released with the code.

\section{Limitations}
\label{sec:limitations}

The study is offline and pairwise. It measures the quality of the similarity decision, not
the dynamics of a live cache, so cache size, eviction, and the time-evolution of the stored
set are out of scope; the false-hit rate we report is a property of the decision rule that
any such system inherits. We do not evaluate exact-match or two-tier caches as separability
baselines, because our pairs are non-identical by construction and an exact-match layer
never fires on them; a two-tier system simply adds our semantic layer on top of a free exact
layer, so our frontier is the part that governs its risk. The downstream cost in
\cref{sec:downstream} comes from controlled fault injection rather than from a production
incident log, and we report it under several injection models to show the conclusion does
not hinge on one. The wrong-passage condition injects an unrelated passage, which is harsher
than the near-duplicate a real false hit would serve, so its cost is best read as an upper
bound; the truncated-passage condition brackets the gentler end, and a real false hit sits
between them.

\section{Broader impact}
\label{sec:impact}

A false hit is a silent error: the system returns a confident answer that no log flags as
wrong. In low-stakes settings this trades a little quality for a large cost saving, which is
the point of caching. In settings where a wrong answer carries real cost, such as health,
legal, or financial guidance, the same mechanism can serve stale or incorrect information
without any signal to the user. Our results argue that such systems should not run a single
global similarity threshold; they should calibrate per domain to a false-hit ceiling tied to
the cost of being wrong, and they should monitor the served stream rather than trust it.
Embedding caches also have a documented adversarial surface: an attacker who can shape
prompts can manufacture collisions that force false hits~\cite{zhang2026keycollision}. Our
study quantifies the benign version of the same failure, and the defense is aligned in both
cases, namely a verification step or a conservative, cost-aware threshold rather than blind
trust in similarity. On the other side of the ledger, caching reduces the compute and energy
of repeated inference, and a more reliable cache lets a system claim that saving without
silently degrading the answers it serves.
